\documentclass[10pt,a4paper]{amsart}
\usepackage[margin=19mm]{geometry}
\usepackage{amsmath,amssymb,mathtools}
\usepackage[scr=rsfs,cal=euler]{mathalfa}
\usepackage{xfrac}
\usepackage[unicode,hidelinks,bookmarks=false]{hyperref}
\newcommand{\ie}{i.e.\ }
\newcommand{\cf}{cf.\ }
\newcommand{\resp}{respectively\ }
\newcommand{\Subsection}{\subsection}
\providecommand{\nobreakdash}{\nobreak\hbox{-}}
\setlength{\emergencystretch}{2em}
\multlinegap0pt
\usepackage{bm}
\usepackage{bbm}
\usepackage{dsfont}
\usepackage{xspace}
\usepackage{cancel}
\usepackage{mathtools}
\usepackage{amsthm}
\makeatletter
\@ifundefined{Theorem}{\newtheorem{Theorem}{Theorem}}{}
\@ifundefined{Proposition}{\newtheorem{Proposition}[Theorem]{Proposition}}{}
\makeatother
\title[On $2$-blocks with quaternion defect groups]{On $2$-blocks with quaternion defect groups}
\author{Charles Eaton, Florian Eisele, Radha Kessar, Markus Linckelmann,, A. A. Schaeffer Fry}
\date{Source: arXiv:2608.12052v1 (CC BY 4.0)}
\begin{document}
\maketitle

\noindent\emph{Source and licence.} On $2$-blocks with quaternion defect groups. Version arXiv:2608.12052v1 (CC BY 4.0), distributed under the Creative Commons Attribution 4.0 International licence, as recorded in the arXiv licence metadata for this exact version and shown on the abstract page https://arxiv.org/abs/2608.12052v1. Full rights notice: licenses/arxiv-2608.12052-CC-BY-4.0.txt.\par\smallskip

\noindent\textit{Editorial scope.} Counted unit: the Proposition whose source label is \texttt{prop:Q2A} in the article (source0.tex lines 694--696, source label \texttt{prop:Q2A}), delivered together with the article's own complete original proof of it (source0.tex lines 697--753, one \texttt{proof} environment of 57 lines). No mathematics is written, repaired or re-derived: statement and proof are the article's own text line for line. Required context retained uncounted: none. Every definition, display and label of those lines is retained unchanged. Omitted from this package and not redistributed: 1--693, 754--1914. Nothing inside a retained excerpt is omitted or altered: every retained body is delivered exactly as the article's lines are written, so no omission receipt of any kind accompanies this package. Citations: the retained excerpts carry no citation command at all, so no bibliography entry of the article is transcribed and no reference is redistributed. Reference adaptation: none was needed and none was made; every reference inside the delivered unit resolves inside it, so no unresolved reference or citation is produced. Printed slips kept literal: no printed word, symbol, hypothesis or number of the counted statement or of its proof was altered. Layout and class: the article's own class is \texttt{amsart} with packages including fontenc, amssymb, enumerate, amsthm, amsmath, xcolor, xy, hyperref, theoremref; the wrapper is the lane's pinned \texttt{amsart} editorial wrapper at 10pt on A4, which restates the theorem-like environments the retained text needs and adds the article's own macro definitions that the retained text needs (none), with extra wrapper packages (bm, bbm, dsfont, xspace, cancel, mathtools). The delivered numbering is the unit's own. Assets: no retained span contains a figure environment, a graphics-inclusion command, a PDF-page inclusion, a table or a TikZ picture; no image file is copied into this package, the package declares no asset dependency and no image file is redistributed with it. Licence: version arXiv:2608.12052v1 of this article is distributed under the Creative Commons Attribution 4.0 International licence, as recorded in the arXiv licence metadata for this exact version and shown on the abstract page https://arxiv.org/abs/2608.12052v1. A full rights notice is delivered at licenses/arxiv-2608.12052-CC-BY-4.0.txt. Characters: every retained excerpt is pure ASCII, so the delivered engine, XeLaTeX with its default Latin Modern text font, reports no missing character.
\begin{Proposition}\label{prop:Q2A}
  Let $n\geq 4$ be a natural number and let $c\in k$ be arbitrary. Then the algebra $Q(2A)^{2^{n-2},c}$ has Morita--Frobenius number $1$ if and only if $Q(2A)^{2^{n-2},c} \cong Q(2A)^{2^{n-2},\tilde c}$ with $\tilde c\in \mathbb F_2$.
\end{Proposition}

\begin{proof}
  We will show that for any $c,\tilde c\in k$ there is an isomorphism $Q(2A)^{2^{n-2},c} \cong Q(2A)^{2^{n-2},\tilde c}$ if and only if either $c=\tilde c=0$ or both $c$ and $\tilde c$ are non-zero and $\frac{\tilde c}{c}$ is a $(5\cdot 2^{n-3}-3)$-rd root of unity. Given any $c\neq 0$, the latter part of the assertion means that $Q(2A)^{2^{n-2},c^2} \cong Q(2A)^{2^{n-2},c}$, i.e. $Q(2A)^{2^{n-2},c}$ has Morita--Frobenius number $1$, if and only if $Q(2A)^{2^{n-2},c} \cong Q(2A)^{2^{n-2},1}$. This clearly implies the claim. 

  We prove the ``only if'' direction first. Set $A=Q(2A)^{2^{n-2},c}$ and $\tilde A=Q(2A)^{2^{n-2},\tilde c}$. Assume that $\varphi:\ \tilde A \longrightarrow A$ is an isomorphism of $k$-algebras. To avoid confusion, we denote the generators of $A$ by $e_0$, $e_1$, $\alpha$, $\beta$ and $\gamma$ and those of $\tilde A$ by $\tilde e_0$, $\tilde e_1$, $\tilde \alpha$, $\tilde \beta$ and $\tilde \gamma$. We can assume without loss of generality that $\varphi$ maps $e_i$ to $\tilde e_i$ for $i\in\{0,1\}$, since the Cartan numbers associated to $e_0$ and $e_1$ are different. 
  We define $s=2^{n-2}$ (Erdmann calls this ``$k$'', but this conflicts with our notation), $z=\alpha\beta\gamma+\beta\gamma\alpha+\gamma\alpha\beta$ and $Z=k[z]$.  One checks that $Z\subseteq Z(A)$. Let $\bar{\phantom{f}}: Z\longrightarrow k$ denote the $k$-algebra homomorphism that maps an element of $Z$ to its constant coefficient. 
  
  The relations of $A$ imply that $A$ has a basis consisting of paths in which an $\alpha$ can only be followed by a $\beta$, a $\beta$ can only be followed by a $\gamma$ and a $\gamma$ can only be followed by an $\alpha$. Any path that does not satisfy this lies in the second socle of $A$, and every path of length strictly greater than $3s$ equals zero in $A$. To be more precise, we will use that
  $e_0Ae_1$ has basis $\{ (\beta\gamma\alpha)^i \beta, (\alpha\beta\gamma)^i \alpha\beta \ | \ i=0,\ldots,s-1\}$, 
   $e_1Ae_0$ has basis $\{ (\gamma\alpha\beta)^i \gamma, (\gamma\alpha\beta)^i \gamma\alpha \ | \ i=0,\ldots,s-1\}$ and 
   $e_0Ae_0$ has basis $\{(\alpha\beta\gamma)^i, (\beta\gamma\alpha)^{i+1}, (\alpha\beta\gamma)^i\alpha, (\beta\gamma\alpha)^i\beta\gamma \ | \ i=0,\ldots,s-1\}$. Note that $(\alpha\beta\gamma)^s=(\beta\gamma\alpha)^s$. Also note that the only non-zero path in $A$ having $\alpha^2$ as a proper subpath is $\alpha^3$.

  Given the above, we can assume that
  $$
    \varphi(\tilde \beta)=b\beta + y_1 \alpha\beta \quad\textrm{with $b\in Z^\times$, $y_1\in Z$},
  $$
  $$
    \varphi(\tilde \gamma)=g\gamma + y_2 \gamma\alpha \quad\textrm{with $g\in Z^\times$, $y_2\in Z$},\textrm{ and}
  $$
  $$
    \varphi(\tilde \alpha)=a\alpha + y_3 \beta\gamma+y_4 \alpha\beta\gamma + y_5 \beta\gamma\alpha \quad\textrm{with $a\in Z^\times$, $y_3,y_4,y_5\in Z$}.
  $$
  We get 
  $$
    \varphi(\tilde\beta)\varphi(\tilde\gamma)\varphi(\tilde\beta) = b^2g \beta\gamma\beta + b(by_2+gy_1)\beta\gamma\alpha\beta+y_1(gy_1+by_2)\alpha\beta\gamma\alpha\beta
  $$
  whereas 
  $$
    \varphi(\tilde\beta\tilde\gamma\tilde\beta)= \varphi((\tilde \alpha\tilde\beta\tilde\gamma)^{s-1}\tilde\alpha\tilde\beta)= a^sb^sg^{s-1} (\alpha\beta\gamma)^{s-1}\alpha\beta.
  $$
  It follows that $by_2+gy_1$ annihilates $\beta\gamma\alpha\beta$, which means $by_2+gy_1\in z^{s-1}Z$. But that means $by_2+gy_1$ annihilates $\alpha\beta\gamma\alpha\beta$ as well. It follows that $\varphi(\tilde\beta)\varphi(\tilde\gamma)\varphi(\tilde\beta) = b^2g \beta\gamma\beta$, and therefore 
  \begin{equation} \label{eqn:rels for hom 1}
    \bar b^2 \bar g= \bar a^s \bar b^s \bar g^{s-1}.
  \end{equation} 
  Now compare
  $$
    \varphi(\tilde\alpha)^2 = a^2\alpha^2 + ay_3 \alpha\beta\gamma + ay_5\alpha\beta\gamma\alpha +ay_3\beta\gamma\alpha + y_3y_4 \beta\gamma\alpha\beta\gamma + ay_4 \alpha\beta\gamma\alpha + y_4^2 (\alpha\beta\gamma)^2 + y_3y_5 \beta\gamma\alpha\beta\gamma + y_5^2 (\beta\gamma\alpha)^2 $$
  and 
  $$ \varphi(\tilde \alpha^2)= b^sg^sa^{s-1}(\beta\gamma\alpha)^{s-1}\beta\gamma+\tilde c b^sg^sa^s(\beta\gamma\alpha)^{s},$$
  where we again used $by_2+gy_1\in z^{s-1}Z$ to obtain the latter equation. It follows that $ay_3 \alpha\beta\gamma$ is a scalar multiple of $(\alpha\beta\gamma)^s$, which forces $y_3\in z^{s-1}Z$. But then $ay_3 \alpha\beta\gamma = ay_3\beta\gamma\alpha$ since $(\alpha\beta\gamma)^s=(\beta\gamma\alpha)^s$. That is, $ay_3 \alpha\beta\gamma + ay_3\beta\gamma\alpha =0$ as $k$ has characteristic $2$. Similarly, we see that $a(y_4+y_5)\alpha\beta\gamma\alpha$ must be zero, which implies $y_4+y_5\in z^{s-1}Z$. But then $y_3(y_4+y_5)\beta\gamma\alpha\beta\gamma=0$ and $(y_4+y_5)^2(\beta\gamma\alpha)^2=0$, so 
  $$\varphi(\tilde\alpha)^2 = a^2\alpha^2 = a^2(\beta\gamma\alpha)^{s-1}\beta\gamma+a^2c(\beta\gamma\alpha)^{s}.$$ 
  This gives us 
  \begin{equation} \label{eqn:rels for hom 2}
    \bar a^2 = \bar b^s \bar g^s\bar a^{s-1} \quad \textrm{and} \quad \bar a^2 c= \tilde c \bar b^s \bar g^s \bar a^s.
  \end{equation}
  The last equation implies that $c$ and $\tilde c$ are either both zero or both non-zero. If they are both zero, then we are done. So assume that $c$ and $\tilde c$ are both non-zero. 
  In that case, equations~\eqref{eqn:rels for hom 1}~and~\eqref{eqn:rels for hom 2} imply 
  $$
    \frac{c}{\tilde c} = \bar a \quad\textrm{and} \quad \bar b \bar g =\bar a^{\frac{3-s}{s}} \quad \textrm{and} \quad \bar a^{\frac{5s-6}{s}}=1.
  $$
  Note that $s$ is a power of two, so $s$-th roots are unique in $k$. In particular, $\bar a$ and therefore $\frac{c}{\tilde c}$ is a $(5s-6)$-th root of unity and therefore also a $(5\cdot 2^{n-3}-3)$-rd root of unity, as claimed.

  For the converse assume that $\frac{c}{\tilde c}$ is a $(5\cdot 2^{n-3}-3)$-rd root of unity (there is nothing to show in the case $c=\tilde c=0$). Define $a=\frac{c}{\tilde c}$, $b=a^{\frac{3-s}{s}}$ and $g=1$. Then one checks that the assignment
  $$
    \tilde \alpha \mapsto a \alpha, \ \tilde \beta \mapsto b \beta, \ \tilde \gamma \mapsto g \gamma
  $$
  extends to a $k$-algebra isomorphism between $\tilde A$ and $A$, proving the claim.
\end{proof}

\end{document}
